\chapter{Ergebnisse}
\label{ch:results}

\begin{center}
\fbox{\parbox{0.95\linewidth}{\small
\textbf{Stand: Versuche noch nicht durchgefuehrt.} Dieses Kapitel legt die
Darstellungsform fest, bevor Messwerte vorliegen. Damit steht vorab fest,
welche Kennwerte berichtet werden, und nicht erst im Nachhinein diejenigen,
die sich als guenstig erweisen. Saemtliche Zellen sind Platzhalter und mit
einem Strich gefuellt. Sie sind aus den Pruefprotokollen zu uebernehmen. In
diesem Kapitel duerfen keine Auswertungen erscheinen, die nicht bereits in
\Cref{ch:strategy} festgelegt sind.}}
\end{center}

\section{Beschreibung der Stichprobe}
\label{sec:descriptives}

Zunaechst ist zu dokumentieren, welche Balken in die Untersuchung eingegangen
sind und welche nach den in \Cref{ch:strategy} festgelegten Kriterien
ausgeschlossen wurden. Der Ausschluss wegen unsicherer Zonenzuordnung ist
dabei gesondert auszuweisen, da er die zentrale methodische Schwachstelle der
Untersuchung betrifft.

\begin{table}[H]
\centering
\caption{Zusammensetzung der Stichprobe.}
\label{tab:sample}
\begin{tabular}{@{}lr@{}}
\toprule
& \textbf{Anzahl} \\
\midrule
Bereitgestellte Balken                        & -- \\
Davon mit gesicherter Einbaulage              & -- \\
Ausgeschlossen: Zonenzuordnung unsicher       & -- \\
Ausgeschlossen: Schaedlingsbefall oder Faeule & -- \\
Ausgeschlossen: sonstige Gruende              & -- \\
\midrule
In die Auswertung eingegangene Balken         & -- \\
Probekoerper Druckzone                        & -- \\
Probekoerper Zugzone                          & -- \\
\bottomrule
\end{tabular}
\end{table}

\section{Rohdichte und Holzfeuchte}
\label{sec:density-moisture}

Rohdichte und Holzfeuchte sind vor der Festigkeitsauswertung zu berichten. Die
Rohdichte geht als Kovariate in die Auswertung ein, weil sie die mechanischen
Eigenschaften staerker bestimmt als jede andere leicht messbare Groesse. Ein
systematischer Rohdichteunterschied zwischen den beiden Zonen waere ein
Befund fuer sich und wuerde die Interpretation eines Festigkeitsunterschieds
grundlegend veraendern.

\begin{table}[H]
\centering
\caption{Rohdichte und Holzfeuchte je Zone.}
\label{tab:density}
\begin{tabular}{@{}lrrrr@{}}
\toprule
& \textbf{n} & \textbf{Mittelwert} & \textbf{Standardabw.}
& \textbf{Min / Max} \\
\midrule
Rohdichte Druckzone     & -- & -- & -- & -- \\
Rohdichte Zugzone       & -- & -- & -- & -- \\
Holzfeuchte Druckzone   & -- & -- & -- & -- \\
Holzfeuchte Zugzone     & -- & -- & -- & -- \\
\bottomrule
\end{tabular}
\end{table}

\section{Festigkeitskennwerte je Zone}
\label{sec:strength-results}

Die Festigkeitskennwerte sind getrennt nach ehemaliger Druck- und Zugzone
auszuweisen. Neben Mittelwert und Streuung ist das Fraktil anzugeben, das fuer
eine Einstufung nach \parencite{en384} massgeblich waere, sofern der
Stichprobenumfang eine solche Angabe traegt. Traegt er sie nicht, ist das
ausdruecklich zu vermerken, anstatt einen Wert zu berichten, der eine
Genauigkeit vortaeuscht.

\begin{table}[H]
\centering
\caption{Festigkeitskennwerte je Zone.}
\label{tab:strength}
\begin{tabular}{@{}lrrrr@{}}
\toprule
& \textbf{n} & \textbf{Mittelwert} & \textbf{Standardabw.}
& \textbf{Min / Max} \\
\midrule
Druckzone & -- & -- & -- & -- \\
Zugzone   & -- & -- & -- & -- \\
\bottomrule
\end{tabular}
\end{table}

\section{Vergleich der beiden Zonen}
\label{sec:zone-comparison}

Der Vergleich erfolgt gepaart, da beide Probekoerper eines Paares aus
demselben Balken stammen. Zu berichten sind die mittlere Differenz, das
zugehoerige Konfidenzintervall, die Pruefgroesse und der p-Wert, ferner das
Ergebnis der Pruefung auf Normalverteilung der Differenzen. Ein nicht
signifikantes Ergebnis ist als Ergebnis zu berichten und nicht als
ausgebliebener Befund.

\begin{table}[H]
\centering
\caption{Gepaarter Vergleich Druckzone gegen Zugzone.}
\label{tab:paired-test}
\begin{tabular}{@{}lr@{}}
\toprule
& \textbf{Wert} \\
\midrule
Anzahl Paare                          & -- \\
Mittlere Differenz                    & -- \\
Konfidenzintervall der Differenz      & -- \\
Pruefgroesse                          & -- \\
p-Wert                                & -- \\
Effektgroesse                         & -- \\
Pruefung auf Normalverteilung         & -- \\
\bottomrule
\end{tabular}
\end{table}

\section{Auswertung unter Beruecksichtigung der Rohdichte}
\label{sec:covariate}

Ergibt sich zwischen den Zonen ein Rohdichteunterschied, ist der
Festigkeitsvergleich zusaetzlich unter Beruecksichtigung der Rohdichte als
Kovariate zu fuehren. Nur so laesst sich trennen, ob ein beobachteter
Unterschied auf die \gls{lastgeschichte} zurueckgeht oder schlicht darauf,
dass die Probekoerper der beiden Zonen aus unterschiedlich dichtem Material
stammen.

\begin{table}[H]
\centering
\caption{Auswertung mit Rohdichte als Kovariate.}
\label{tab:ancova}
\begin{tabular}{@{}lrrr@{}}
\toprule
\textbf{Quelle} & \textbf{Freiheitsgrade} & \textbf{Pruefgroesse}
& \textbf{p-Wert} \\
\midrule
Zone                     & -- & -- & -- \\
Rohdichte                & -- & -- & -- \\
Balken (Blockeffekt)     & -- & -- & -- \\
\bottomrule
\end{tabular}
\end{table}
